\documentclass[12pt]{article}
\usepackage[utf8]{inputenc}
\usepackage[T1]{fontenc}
\usepackage{hyperref}
\title{Shanenawa Language Resources for K-12 Curriculum Development}
\author{Compiled by Claude Sonnet 5.5 (claude-sonnet-5-5) for Andr\'as Kornai}
\date{}
\begin{document}
\maketitle

\section{Introduction}
Shanenawa (ISO 639-3: swo; Glottocode shan1283) is an endangered Panoan language of Acre, Brazil, spoken by about nine people, mostly elderly. Search results place the community along the Envira River in the municipality of Feij\'o. Younger people increasingly use Portuguese. The searches found only academic descriptions, all in Portuguese. They found no school materials, primers or dictionaries.

\section{Primary Educational Resources}
The searches found no primers, textbooks or school programmes for Shanenawa. Any curriculum would have to be built from the linguistic descriptions below, with the community.

\section{Linguistic and Reference Materials}
\begin{itemize}
\item Morfossintaxe da l\'ingua Shanen\'awa (Pano): as classes de palavras e sua estrutura interna (doctoral work, Universidade Federal de Goi\'as repository): \url{https://repositorio.bc.ufg.br/items/10533295-6d16-44f4-99f2-a3605ccfc50b}
\item Sintagmas nominais plenos: a marca de transitividade na l\'ingua Shanenawa (Fam\'ilia Pano), Universidade Federal de Goi\'as repository: \url{https://repositorio.bc.ufg.br/items/44034b29-8d23-48db-99b1-d7306a72b529}
\item Article in \emph{Via Litterae} on Shanenawa morphology (Universidade Estadual de Goi\'as): \url{https://www.revista.ueg.br/index.php/vialitterae/article/download/5339/3600/18449}
\item Aspectos fonol\'ogicos da l\'ingua Shanenawa (Pano), UNICAMP (the search summary says it includes a basic vocabulary appendix; I did not open it): \url{https://www.repositorio.unicamp.br/acervo/detalhe/135131}
\item Aspectos de uma gram\'atica Shaw\~a (Pano), UNICAMP. Shaw\~a is the name used for a related Panoan language, so check this entry before relying on it: \url{https://www.repositorio.unicamp.br/acervo/detalhe/871272}
\item Uma proposta de classifica\c{c}\~ao interna das l\'inguas da fam\'ilia pano (Panoan classification, mentions Shanenawa): \url{https://periodicos.ufpe.br/revistas/INV/article/view/1443}
\end{itemize}

\section{Religious and Community Resources}
The searches found no scripture translations or community organisations for Shanenawa.

\section{Digital and Audio Resources}
The searches found no audio, video or software resources specific to Shanenawa.

\section{Curriculum Development Context}
Almost nothing exists for classroom use. The first steps would be documentation with the remaining elderly speakers, an agreed orthography, and a vocabulary built from the academic descriptions above. Work should be led by the community. The sources are in Portuguese.
\end{document}